\section{Introduction}

Public blockchain networks like Ethereum are funda-\
mentally constrained by the blockchain trilemma, limit-\
ing their Layer-1 (L1) base-layer throughput to 15-25\
transactions per second (TPS) and forcing the network to\
deal with congestion and volatile transaction fees .
To scale Ethereum’s ecosystem, rollups have been adopted\
as the main roadmap for scaling transactions. Layer-2\
(L2) ZK-Rollups off-chain execute transactions, compress\
their data and validity proofs (SNARKs/STARKs) and submit\
them to the Ethereum blockchain for verification via\
smart contract. The Dencun upgrade on Ethereum\
implemented EIP-4844, known as Proto-Danksharding,\
which introduced blob-carrying transactions that\
reduced the data costs for Layer-2 rollups by over 90%,\
fueling a significant increase in Ethereum transactions~\cite{eip4844, l2beat_dencun, dune_eip4844}.

However, the multi-dimensional parameter space of rollups---comprising batch sizes (\(B_{\text{max}}\)), timeouts (\(T_{\text{timeout}}\)), sequencing rules, DA modes, and provers remains poorly understood in controlled settings. In production rollups (e.g., ZKsync Era~\cite{zksync_era_docs}, Starknet~\cite{starknet_docs}, Scroll~\cite{scroll_docs}), these components are tightly coupled with complex operational mechanisms, making isolated scientific analysis of individual configurations impossible. To address this, we present \textbf{RollupX}, a modular L2 ZK-rollup benchmarking sandbox implemented in Rust and Solidity that leverages a Sparse Merkle Tree (SMT) state manager, the RISC0 zkVM for execution/proving, and a Groth16 wrapper for proof compression.

\subsection{Rationale for a Clean-Room Sandbox}

The justification for building a clean-room, modular ZK-rollup sandbox like RollupX rather than simply forking or benchmarking an existing production ZK-rollup (such as ZKsync Era~\cite{zksync_era_docs}, Starknet~\cite{starknet_docs}, or Scroll~\cite{scroll_docs})---is grounded in system isolation, architectural agility, and research reproducibility:

\begin{enumerate}

\item \textbf{State Representation:} Ethereum uses a Hexary Patricia Trie (PMT) to manage state~\cite{wood2014ethereum}, which is very slow to prove in ZK due to variable path lengths and complex serialization~\cite{trilemma}. Similar to how production rollups use binary tries to minimize ZK circuit complexity and prover latency~\cite{zksync_era_docs, starknet_docs}, RollupX implements a 256-level binary Sparse Merkle Tree (SMT) to guarantee simple, constant-length ZK proofs.

\item \textbf{zkVM Execution:} Writing custom ZK circuits (like in Circom) is complex, hard to audit, and inflexible. RollupX compiles standard Rust code directly inside the RISC0 zkVM, making the state transition logic easy to audit and modify~\cite{lavaur2023modular}. While zkVM execution has instruction overhead, parallel proving and recursive merging make it commercially viable for L2 protocols~\cite{werner2024analyzing}.

\item \textbf{EIP-4844 Blob Dilemma:} Under EIP-4844, blobs are sold in fixed 128~KB chunks~\cite{eip4844}. For low-volume rollups, this creates a ``Blob Dilemma'': wait to fill the blob (increasing delay) or post it partially empty (increasing cost)~\cite{werner2024analyzing}. RollupX provides a sandbox to sweep workloads and test how sequencers can dynamically adjust batch parameters to close this cost gap~\cite{adler2021data, l2beat_dencun, dune_eip4844}.

\end{enumerate}

\subsection{Contributions}
Our contributions are threefold:
\begin{enumerate}
    \item \textbf{Empirical Characterization:} We sweep batch sizes, timeouts, and workloads to disclose the complete decoupling of the complexity of executing L2 virtual machine code from the cost of L1 gas.
    \item \textbf{Data Availability \& Prover Analysis:}We compare the cost of Calldata, Blob and Offchain DA modes, and model the RISC0 prover execution to reveal that proving time proceeds in a step-function of the segment boundaries.
    \item \textbf{Open-Source Artifacts:} We release the modular rollup code, Docker container setups, benchmark scripts, telemetry logs and analysis notebooks to allow others to explore and verify our findings.
\end{enumerate}